\chapter{Loops}%
\label{ch:loops}

Computers are tireless. Where a person would balk at adding the numbers from 1 to
a million, a computer does it without complaint \emph{if} we can express the
work as a repetition. A \textbf{loop} runs a block of code over and over, and it
is what turns a few lines of source into millions of actions. Salam gives you
three looping tools: \code{until}, the counting \code{repeat}, and the
collection-walking \code{each}. We will meet all three.

\section{The \texttt{until} loop}

An \code{until} loop repeats its body \emph{as long as} a condition stays true. It
checks the condition, runs the body if true, then checks again round and round
until the condition becomes false.

\begin{salam}
func main:
    mut i : int = 1
    until i <= 3:
        print i
        print " "
        i += 1
    end
    println ""
end
\end{salam}

\begin{progout}
1 2 3
\end{progout}

Trace it: \code{i} starts at 1; \code{1 <= 3} is true, so we print \code{1} and
bump \code{i} to 2; and so on until \code{i} becomes 4, at which point
\code{4 <= 3} is false and the loop stops. The variable that drives the loop ---
here \code{i} must be \code{mut}, because the loop changes it.

\begin{warn}
Every \code{until} loop needs a way to \emph{end}. If the condition never becomes
false for example if you forget the \code{i += 1} above the loop runs
forever and your program hangs. Whenever you write an \code{until}, find the line
that makes progress toward the exit and check that it is really there.
\end{warn}

\code{until} shines when you do not know in advance how many times you will loop.
A classic example is Euclid's algorithm for the greatest common divisor, which
loops until a remainder reaches zero:

\begin{salam}
func gcd(a : int, b : int) : int:
    mut x : int = a
    mut y : int = b
    until y != 0:
        mut t : int = y
        y = x % y
        x = t
    end
    ret x
end

func main:
    println gcd(48, 36)
end
\end{salam}

\begin{progout}
12
\end{progout}

\section{The \texttt{repeat} loop}

When you \emph{do} know the pattern of counting a set number of times, a range
of numbers, or every $S$-th number in that range \code{repeat} is the tool.
It reads almost like plain English, and it is Salam's direct replacement for the
classic counting loop found in other languages. It comes in a few forms.

\subsection{Repeat a fixed number of times}

\code{repeat N:} runs its body \code{N} times. There is no loop variable to
manage when you only need the \emph{repetition}, this is the cleanest tool:

\begin{salam}
func main:
    repeat 3:
        println "Salam!"
    end
end
\end{salam}

\begin{progout}
Salam!
Salam!
Salam!
\end{progout}

If the count is zero or negative, the body simply does not run.

Sometimes you also want to know \emph{which} pass you are on, counting from
zero. Add \code{with} and a name to bind it:

\begin{salam}
func main:
    repeat 5 with i:
        print i
    end
    println ""
end
\end{salam}

\begin{progout}
01234
\end{progout}

\code{repeat 5 with i} binds \code{i} to \code{0}, \code{1}, \code{2}, \code{3},
then \code{4} one pass per value. This is the direct replacement for the
classic ``count from 0 up to $n - 1$'' loop that other languages spell with a
three-part \code{for} header.

\begin{warn}
Do not spell ``count from 0 up to $n - 1$'' as \code{repeat 0 to n - 1 with i}.
It looks reasonable, but when \code{n} is \code{0} the range \code{0 to -1} is
read as counting \emph{downward}, so the body still runs once when you meant it
to run zero times. The safe, idiomatic way to write that pattern is
\code{repeat n with i}, exactly as above; let \code{repeat} pick the count,
do not hand it a hand-built range that can point the wrong way.
\end{warn}

\subsection{Repeat over a range}

\code{repeat A to B:} runs the body once for every value from \code{A} up to
\code{B} \emph{inclusive} so \code{repeat 1 to 5} runs five times. When you
just need the repetition and not the value, that is all you need:

\begin{salam}
func main:
    mut sum : int = 0
    mut k : int = 0
    repeat 1 to 5:
        k += 1
        sum += k
    end
    println "1+2+3+4+5 =", sum
end
\end{salam}

\begin{progout}
1+2+3+4+5 = 15
\end{progout}

More often, though, you want to read the current value directly instead of
tracking your own counter. Add \code{with} and a name, just like with
\code{repeat N}:

\begin{salam}
func main:
    repeat 1 to 10 with n:
        println "7 x", n, "=", 7 * n
    end
end
\end{salam}

\begin{progout}
7 x 1 = 7
7 x 2 = 14
...
7 x 10 = 70
\end{progout}

\begin{note}
The upper bound is \emph{inclusive}: \code{repeat 1 to 5} covers
1,~2,~3,~4,~5 five iterations, not four. This matches how we count in everyday
speech (``count from 1 to 5''). Keep the distinction in mind whenever you set
up a range: the loop always includes both endpoints.
\end{note}

\subsection{Repeat with a step}

Add \code{by S} to stride by \code{S} instead of 1. Combine it with \code{with}
to read each value as you go:

\begin{salam}
func main:
    repeat 0 to 20 by 2 with i:
        print i
        print " "
    end
    println ""
end
\end{salam}

\begin{progout}
0 2 4 6 8 10 12 14 16 18 20
\end{progout}

If you only care how many times the loop ran and not the values themselves,
drop \code{with} just as with any other \code{repeat} range:

\begin{salam}
func main:
    mut count : int = 0
    repeat 0 to 20 by 2:
        count += 1
    end
    println "even numbers in 0..20:", count
end
\end{salam}

\begin{progout}
even numbers in 0..20: 11
\end{progout}

There are eleven of them (0,~2,~\dots,~20), and the loop counts each another
reminder that the upper bound is included.

\begin{tip}
\code{repeat} covers counting, but it is not the tool for walking the elements
of an array or other collection directly. That job belongs to \code{each}, which
we meet properly in Chapter~7: \code{each x in numbers:} hands you each element
in turn, with no index bookkeeping at all.
\end{tip}

\section{Breaking out and skipping ahead}

Two keywords give you finer control \emph{inside} any loop:

\begin{itemize}[leftmargin=1.4em]
  \item \code{break} stops the loop immediately and continues with whatever comes
        after it.
  \item \code{continue} abandons the current pass and jumps straight to the next
        one.
\end{itemize}

\begin{salam}
func main:
    repeat 6 with n:
        if n == 3: continue end      // skip 3
        if n == 5: break end         // stop before 5
        print n
    end
    println ""
end
\end{salam}

\begin{progout}
0124
\end{progout}

The loop prints 0, 1, and 2 normally; at \code{n == 3} the \code{continue} skips
the \code{print}, so 3 never appears; at \code{n == 5} the \code{break} ends the
loop before printing. Hence \code{0124}.

\begin{tip}
\code{break} and \code{continue} are handy, but a loop peppered with them can be
hard to follow. Often a clearer condition in the \code{until} or \code{repeat}
header expresses the same intent without jumping. Reserve \code{break} for cases
like ``search until found,'' where stopping early is the whole point.
\end{tip}

\section{Nested loops}

A loop body may itself contain a loop. Nested loops are how you work with grids,
tables, and shapes. Here a loop over rows contains a loop over columns to draw a
triangle of stars using the ``print then end the line'' trick from Chapter~1:

\begin{salam}
func main:
    mut row : int = 0
    repeat 1 to 5:
        row += 1
        repeat 1 to row:
            print "*"
        end
        println ""
    end
end
\end{salam}

\begin{progout}
*
**
***
****
*****
\end{progout}

The outer loop runs once per row; for each row, the inner loop prints as many
stars as the row number, and then \code{println ""} moves to the next line.
Nested loops multiply: an outer loop of 5 around an inner loop of up to 5 does up
to 25 iterations in total.

\begin{warn}
\code{break} and \code{continue} affect only the \emph{innermost} loop that
encloses them. A \code{break} inside the inner star loop would stop that row's
stars, not the whole triangle.
\end{warn}

\section{In Persian}

The loop keywords have Persian forms too: \code{\fa{تاوقتی}} is the Persian
spelling of \code{until}, \code{\fa{تکرار}} is the Persian spelling of
\code{repeat}, and \code{\fa{هر}} is the Persian spelling of \code{each}. Inside
a \code{repeat}, \code{\fa{تا}} spells the range word \code{to}, and
\code{\fa{گام}} spells the step word \code{by}. Salam has never had a C-style
\code{for}, so there is no Persian word tied to one \code{repeat} and
\code{each} cover that ground instead. Here is the star loop in Persian:

\begin{salamfa}
تابع اصلی:
    تکرار 1 تا 5:
        بنویس "*"
    تمام
    چاپ ""
تمام
\end{salamfa}

\section{Which loop should I use?}

\begin{itemize}[leftmargin=1.4em]
  \item Reach for \code{repeat N} when you just need to do something a set number
        of times, and \code{repeat A to B} (optionally with \code{by}) when you
        are counting over a clear range.
  \item Add \code{with i} to a \code{repeat} when you need to read the current
        number inside the body especially the classic ``count from 0 up to
        $n - 1$'' pattern, which is simply \code{repeat n with i}.
  \item Use \code{each} when you are walking the elements of an array or other
        collection directly, rather than counting (see Chapter~7).
  \item Use \code{until} when the number of iterations is not known up front and
        the loop ends when some condition changes.
\end{itemize}

\section{Chapter summary}

\begin{itemize}[leftmargin=1.4em]
  \item \code{until cond:} repeats while a boolean condition holds; make sure it
        can end.
  \item \code{repeat N}, \code{repeat N with i}, \code{repeat A to B}, and
        \code{repeat A to B by S} cover fixed counts, indices, and ranges; the
        upper bound of a range is \emph{inclusive}. Never spell ``count from 0''
        as \code{repeat 0 to n - 1 with i} use \code{repeat n with i}.
  \item \code{each x in collection:} walks the elements of a collection directly
        (Chapter~7).
  \item \code{break} ends the innermost loop; \code{continue} skips to its next
        pass.
  \item Loops nest, and that is how you handle grids and shapes.
\end{itemize}

\section{Exercises}

\exercise Use an \code{until} loop to print the numbers from 10 down to 1, then
print \code{"Lift-off!"}.

\exercise Use a \code{repeat} loop to compute and print the sum of all numbers
from 1 to 100.

\exercise Using \code{repeat ... by}, print every multiple of 5 from 5 to 50 on
one line, separated by spaces.

\exercise Use a loop with \code{continue} to print the numbers from 1 to 20,
skipping every multiple of 3.

\exercise Draw an upside-down triangle of stars (five stars on the first row, one
on the last) using nested loops.

\exercise Rewrite the times-table example so it prints the table for a number
stored in a variable, in Persian.
